\section{Background and Related Work}\label{sec:background}

\paragraph{Open agent discovery.}
Project NANDA's current official title is \emph{Architecting the Internet of AI Agents}
\citep{projectnanda2026}.  Its research anchor, ``Beyond DNS: Unlocking the
Internet of AI Agents via the NANDA Index and Verified AgentFacts,'' proposes
indexing and verifiable facts for networked agents
\citep{raskar2025beyonddns}.  This work motivates an open discovery layer, but
\Argorix does not implement NANDA, DNS resolution, or a remote index.  The
comparison is architectural: locally compiled sovereign metadata could become
an input to a future discovery service.

\paragraph{Names, identifiers, and credentials.}
DNS separates hierarchical naming from the mechanisms that applications build
above it \citep{mockapetris1987dnsconcepts,mockapetris1987dnsimplementation}.
DID Core specifies identifiers and DID documents, while the Verifiable
Credentials Data Model specifies a credential data model
\citep{w3c_did_core_2022,w3c_vc_data_model_2025}.  Declaring a string shaped
like a DID or a credential reference is not equivalent to resolving,
authenticating, issuing, or verifying it.  We use that distinction throughout.

\paragraph{Agent interaction protocols.}
MCP defines a protocol for context and tool interactions; A2A defines
inter-agent communication \citep{mcp_spec_2025,a2a_spec}.  \Argorix models
bindings and bridge constraints for these boundaries declaratively.  The
observed runtime records that bridge runtime, network, tool, or agent execution
is disabled.  It therefore tests declarations and negative controls rather
than interoperability with live MCP or A2A endpoints.

\paragraph{Risk framing.}
NIST AI RMF organizes AI risk management as a lifecycle activity
\citep{tabassi2023airmf}, and OWASP catalogs application risks including prompt
injection and sensitive-information disclosure \citep{owasp_llm_top10}.  These
sources inform threat categories, not certification mappings.  A policy name
or governance record cannot establish legal compliance.

\paragraph{Policy as code and agent languages.}
Open Policy Agent (OPA) separates a general-purpose policy decision engine from
application enforcement and expresses decisions in Rego
\citep{opa_docs}.  \Argorix instead compiles policy together with agent,
passport, provider, and evidence declarations and emits VM-local decision
events; it is not an OPA implementation or Rego evaluator.  Jason interprets an
extension of AgentSpeak for belief--desire--intention multi-agent programming
\citep{bordini2007jason}.  Jason's primary concern is agent reasoning and
multi-agent execution, whereas this study examines compiled governance,
provider denial, and runtime evidence.  \Argorix does not claim Jason
compatibility or equivalent BDI semantics.

\paragraph{Provenance and isolation.}
in-toto cryptographically records and verifies software supply-chain steps
\citep{torresarias2019intoto}.  \Argorix EvidenceBundles concern one local
runtime result, use unsigned semantic digests, and exclude source; they neither
implement in-toto layouts nor provide its signer-based supply-chain guarantees.
WASI specifies standards-track host APIs and a capability-based sandbox in
which authority is explicitly granted \citep{wasi_intro}.  The current
\Argorix profile denies operations at the VM/application layer but does not run
inside, implement, or empirically compare a WASI sandbox.

\begin{table}[t]
\centering
\caption{Scope of related-work comparisons. Each row is relational, not an
equivalence or conformance claim.}
\label{tab:related-work}
\begin{tabularx}{\linewidth}{@{}XXX@{}}
\toprule
System & Comparison dimension & Bibliographic scope \\
\midrule
Project NANDA / AgentFacts / ANS & Agent naming and discovery & Project-level comparison only \\
ATrust & Agent trust relationships & Project-level comparison only \\
DCP-AI & Discovery and control plane & Project-level comparison only \\
OAuth/OIDC, FIDO, AP2, Verifiable Intent & Delegation and agent authorization & Conceptual comparison only \\
SPIFFE/SPIRE and OWASP NHI & Non-human/workload identity & Scope comparison only \\
Open Policy Agent & External policy-as-code engine & Direct architectural comparison \\
Cedar & Authorization policy language & Conceptual comparison only \\
Jason / AgentSpeak & BDI agent language and interpreter & Direct language/runtime comparison \\
in-toto / SLSA / VET Your Agent & Evidence, provenance, and external evaluation & Scope comparison only \\
WASI & Capability-based host isolation & Direct isolation-boundary comparison \\
\bottomrule
\end{tabularx}

\end{table}

\Cref{tab:related-work} pairs each system with one comparison dimension.  Its
rows are exact side-by-side comparisons, while the surrounding figure program
shows architectural topology; neither claims implementation equivalence,
protocol conformance, or shared assurance.

\begin{figure}[t]
\centering
\includegraphics[width=\linewidth]{evolution-timeline.pdf}
\caption{\Argorix language evolution represented by repository versions and
feature boundaries. The timeline establishes implementation sequence, not
maturity or deployment readiness.}
\label{fig:evolution}
\end{figure}

The language is implemented in Rust \citep{rust_project}.  The evolution in
\cref{fig:evolution} reflects a sequence from agent execution toward typed
messages, provider contracts, passports, ATrust metadata, governance,
hardening, and evidence.  Version numbers establish repository history only;
they are not external certifications.
